\section*{第 \ref{ch:b3:exoplanets} 章——系外行星}

\begin{solution}{exo:b3:exoplanets:1}
(a)~$K = (2\pi G/P)^{1/3}m_p/\bigl((M_\odot + m_p)^{2/3}\sqrt{1 - e^2}\bigr)$。
木星：$P = \qty{3.742e8}{s}$，$(2\pi G/P)^{1/3} = \num{1.039e-6}$，以国际
单位制计，$K = \num{1.039e-6}\times\num{1.898e27}/(\num{1.990e30})^{2/3}/0.9988
= \qty{12.5}{m.s^{-1}}$。地球：$P = \qty{3.155e7}{s}$，$K =
\qty{0.089}{m.s^{-1}}$，约 \qty{9}{cm.s^{-1}}。(b)~\qty{3.5}{km/h} $=
\qty{0.97}{m.s^{-1}}$，和 \qty{0.35}{km/h} $= \qty{0.097}{m.s^{-1}}$。
木星的信号是第一级精度的十三倍；地球的
信号甚至低于第二级精度，恒星自身的噪声还要叠加在上面。
\end{solution}

\begin{solution}{exo:b3:exoplanets:2}
(a)~$(71492/695700)^2 = \num{0.0106}$，\omterm{def:b3:exoplanets:transit}{凌星深度}为 \qty{1.06}{\percent}；
$(6378/695700)^2 = \num{8.4e-5}$，即百万分之 84。(b)~$R_\star/a$：
地球 $695700/\num{1.496e8} = \num{4.7e-3}$，约 215 分之一的概率；
木星 $695700/\num{7.785e8} = \num{8.9e-4}$，1100 分之一的概率。
\end{solution}

\begin{solution}{exo:b3:exoplanets:3}
$m_p\sin i \approx K M_\odot^{2/3}(P/2\pi G)^{1/3}$，其中 $P =
\qty{3.655e5}{s}$：$(P/2\pi G)^{1/3} = \num{9.55e4}$，以国际单位制计，因此 $m_p
\sin i = 59\times(\num{1.9884e30})^{2/3}\times\num{9.55e4} =
\qty{8.9e26}{kg}$，即 0.47 个木星质量或 149 个地球质量（
对 $m_p$ 在 $M_\star + m_p$ 中的修正不足千分之一）。
$a = \bigl(GM_\odot P^2/4\pi^2\bigr)^{1/3} = \qty{7.66e9}{m} =
\qty{0.051}{au}$。
\end{solution}

\begin{solution}{exo:b3:exoplanets:4}
(a)~$4GM/c^2 = 4\times\num{6.674e-11}\times\num{9.94e29}/(\num{2.998e8})^2 =
\qty{2953}{m}$；$(D_S - D_L)/(D_LD_S) = 1/D_S\times(D_S - D_L)/D_L =
1/(\qty{8}{kpc}) = 1/\qty{2.469e20}{m}$。$\theta_\mathrm E =
\sqrt{2953/\num{2.469e20}} = \qty{3.46e-9}{rad} = \qty{0.71}{mas}$。在
透镜处，$R_\mathrm E = \theta_\mathrm E D_L = \num{3.46e-9}\times
\qty{1.234e20}{m} = \qty{4.27e11}{m} = \qty{2.85}{au}$。(b)~$A(1) =
3/\sqrt5 = 1.34$；$A(0.5) = 2.25/(0.5\times2.062) = 2.18$；$A(0.1) =
2.01/(0.1\times2.002) = 10.0$。(c)~$t_\mathrm E = \num{4.27e11}/\num{2e5}
= \qty{2.13e6}{s}$，约 25 天：\omterm{def:b3:exoplanets:microlensing}{引力微透镜}事件持续数周。
\end{solution}

\begin{solution}{exo:b3:exoplanets:5}
(a)~$a = (GM_\odot P^2/4\pi^2)^{1/3} = \qty{6.09e9}{m}$（\qty{0.0407}{au}），
因此 $a/R_\star = 8.75$。(b)~当 $P/\pi = \qty{22.9}{h}$，对于 $b = 0$：
\[
T_{14} = \frac P\pi\arcsin\frac{1.12}{8.75} = \qty{2.94}{h}, \qquad
T_{23} = \frac P\pi\arcsin\frac{0.88}{8.75} = \qty{2.31}{h}.
\]
对于 $b = 0.7$：$\cos i = 0.7/8.75$，$i = \ang{85.4}$，$(a/R_\star)\sin i =
8.726$，并有
\[
T_{14} = \frac P\pi\arcsin\frac{\sqrt{1.2544 - 0.49}}{8.726} = \qty{2.30}{h},
\qquad
T_{23} = \frac P\pi\arcsin\frac{\sqrt{0.7744 - 0.49}}{8.726} = \qty{1.40}{h}.
\]
(c)~入凌与出凌占凌星时长的 $(T_{14} - T_{23})/T_{14} = \qty{21}{\percent}$
（居中时），$b = 0.7$ 时为 \qty{39}{\percent}；擦
边凌星（$b \to 1$）全无平底，形似字母 V。
\end{solution}

\begin{solution}{exo:b3:exoplanets:6}
$\bar\rho = M/\tfrac43\pi R^3$，其中 $\tfrac43\pi R_\oplus^3 =
\qty{1.087e21}{m^3}$：TRAPPIST-1~b $1.374\times\num{5.97e24}/(1.116^3\times
\num{1.087e21}) = \qty{5.43}{g.cm^{-3}}$；Kepler-10~b
\qty{5.60}{g.cm^{-3}}；WASP-39~b \qty{0.167}{g.cm^{-3}}；HD~209458~b
\qty{0.337}{g.cm^{-3}}。前两者的密度与地球相当
（采用相同常数时为 \qty{5.49}{g.cm^{-3}}），且为岩质；后两者
均为气态巨行星，密度低于水，WASP-39~b 的密度比
木星低八倍。
\end{solution}

\begin{solution}{exo:b3:exoplanets:7}
两极值分别为 $\gamma + K(1 + e\cos\omega)$ 与 $\gamma - K(1 -
e\cos\omega)$。二者之差为 $2K = \qty{60.0}{m.s^{-1}}$，因此 $K =
\qty{30.0}{m.s^{-1}}$；二者的均值为 $\gamma + Ke\cos\omega =
\qty{20.4}{m.s^{-1}}$，因此 $Ke\cos\omega = \qty{8.0}{m.s^{-1}}$，且
$e\cos\omega = 0.27$。在偏心轨道上，恒星在
远星点附近停留，并迅速掠过近星点，因此曲线不对称，
两极值的中点并不等于时间平均值 $\gamma$。
\end{solution}

\begin{solution}{exo:b3:exoplanets:8}
(a)~$g = GM/R^2 = \num{6.674e-11}\times\num{5.332e26}/(\num{9.144e7})^2 =
\qty{4.26}{m.s^{-2}}$；$H = kT/\bar m g = \num{1.381e-23}\times1100/(2.3
\times\num{1.661e-27}\times4.26) = \qty{934}{km}$。取 $R_\star =
\qty{6.533e8}{m}$，\omterm{def:b3:exoplanets:transit}{凌星深度}为 $(\num{9.144e7}/\num{6.533e8})^2 =
\num{0.0196}$。(b)~$2HR_p/R_\star^2 = 2\times\num{9.34e5}\times
\num{9.144e7}/(\num{6.533e8})^2 = \num{4.0e-4}$，即百万分之 400；
两个标高即百万分之 800：低重力的气态巨行星是
\omterm{def:b3:exoplanets:transmission}{透射光谱法}最易观测的目标。
\end{solution}

\begin{solution}{exo:b3:exoplanets:9}
$p = 1.116\times6378/(0.1192\times695700) = 0.0858$，$p^2 = \num{7.37e-3}$。
$B_\lambda \propto 1/(e^{hc/\lambda kT} - 1)$，其中 $hc/\lambda k =
\qty{959}{K}$，在 \qty{15}{\micro m} 处，恒星为 $1/(e^{0.374} - 1) = 2.20$；
日侧在 \qty{400}{K} 时，$1/(e^{2.40} - 1) = 0.0998$，比值为 \num{0.0453}，
\omterm{def:b3:exoplanets:transit}{凌星深度}为 \num{3.3e-4}（即百万分之 334）；在 \qty{500}{K} 时，$1/(e^{1.92} -
1) = 0.172$，比值为 \num{0.078}，\omterm{def:b3:exoplanets:transit}{凌星深度}为 \num{5.7e-4}。温度相差一百开尔文
会使\omterm{def:b3:exoplanets:transit}{凌星深度}变化 \qty{70}{\percent}：实测\omterm{def:b3:exoplanets:transit}{凌星深度}可较精确地确定
日侧温度，较热的日侧意味着
被输送至夜侧的热量很少。
\end{solution}

\begin{solution}{exo:b3:exoplanets:10}
(a)~$S \propto L/d^2$，因此 $d = \sqrt{0.005}\,\unit{au} = \qty{0.0707}{au}$；
$P = 2\pi\sqrt{a^3/GM} = \qty{1.33e6}{s} = \qty{15.4}{d}$。
(b)~$K = (2\pi G/P)^{1/3}M_\oplus/(0.2M_\odot)^{2/3} = \qty{0.75}{m.s^{-1}}$，
是地球对太阳效应的 8.4 倍：低质量恒星运动更明显，且
短周期也有帮助。(c)~$(6378/139140)^2 = \num{2.1e-3}$，是
地球凌日时的 25 倍。两个信号都更大，凌星
每两周出现一次，且凌星概率 $R_\star/a =
\num{0.013}$ 是地球的三倍：小恒星使
宜居带内的岩质行星有机会被发现和研究。
\end{solution}

\begin{solution}{exo:b3:exoplanets:11}
利用 $r\,dr = -\mu\,d\mu$，平均强度为 $\frac1\pi\int_0^1
I\,2\pi r\,dr = 2\int_0^1 I(\mu)\mu\,d\mu = 2\bigl[\tfrac12 - u_1(\tfrac12 -
\tfrac13) - u_2\int_0^1(1 - \mu)^2\mu\,d\mu\bigr]$。末项积分为
$\tfrac12 - \tfrac23 + \tfrac14 = \tfrac1{12}$，故平均值为 $1 - u_1/3 -
u_2/6$，此处为 $0.825$。小型行星遮挡 $\pi p^2 I(\mu)$，相对于 $\pi\times
0.825$，深度为 $p^2 I(\mu)/0.825$。在 $r = 0$ 处，$I = 1$：$1.212\,p^2$；在
$r = 0.5$（$\mu = 0.866$），$I = 0.942$：$1.142\,p^2$；在 $r = 0.9$（$\mu =
0.436$），$I = 0.695$：$0.842\,p^2$。若采用均匀圆盘模型，则会把
$1.212\,p^2$ 当作 $p^2$，从而将半径高估 $\sqrt{1.212} = 1.10$
倍，即 \qty{10}{\percent}。
\end{solution}

\begin{solution}{exo:b3:exoplanets:12}
(a)~$R_\odot/a$：\qty{0.5}{au} 处为 \num{9.3e-3}，\qty{0.7}{au} 处为
\num{6.6e-3}，\qty{1}{au} 处为 \num{4.7e-3}。(b)~$42000\times0.11\times
\num{6.6e-3} = 31$ 颗行星凌日。(c)~其中仅约三分之一被发现：
其余的凌日信号过浅（对地球而言为百万分之 84
）或过少（每年一两次），以致不能高出噪声。这
一比例即探测效率。若不对此及凌日概率加以修正，
则计数会低估此类行星的
普遍程度约 460 倍。
\end{solution}

\begin{solution}{pb:b3:exoplanets:1}
\textbf{1.} $M_\star = 0.80\times\num{1.9884e30} = \qty{1.591e30}{kg}$，$P =
\qty{1.0368e6}{s}$：$a = (GM_\star P^2/4\pi^2)^{1/3} = \qty{1.4245e10}{m}
= \qty{0.0952}{au}$。

\textbf{2.} $v = 2\pi a/P = \qty{86.3}{km.s^{-1}}$。

\textbf{3.} $m_p\sin i = KM_\star^{2/3}(P/2\pi G)^{1/3} = 2.59\times\num{1.363e20}
\times\num{1.352e5} = \qty{4.77e25}{kg}$，\num{7.99} 个地球质量。

\textbf{4.} $a_\star = a\,m_p/M_\star = \num{1.4245e10}\times\num{4.77e25}/
\num{1.591e30} = \qty{427}{km}$；$2\pi a_\star/P = \qty{2.59}{m.s^{-1}}$，
即观测的 $K$。

\textbf{5.} 相对于精度 \qty{1}{m.s^{-1}} 的 $2.59$ 倍，相对于
\qty{0.1}{m.s^{-1}} 的 26 倍：该行星在最佳仪器的探测范围之内，
而上一代仪器则仅勉强可及。

\textbf{6.} $m_p\sin i \propto K\sqrt{1 - e^2}$：$\times0.980$，7.83 个地球
质量，减少 \qty{2}{\percent}。

\textbf{7.} $p = \sqrt{\num{8.38e-4}} = 0.02895$。

\textbf{8.} $R_\star = 0.76\times\qty{695700}{km} = \qty{528700}{km}$；$R_p =
\qty{15310}{km} = 2.40$ 个地球半径。

\textbf{9.} $a/R_\star = \num{1.4245e10}/\num{5.287e8} = 26.94$。

\textbf{10.} $T_{14} = (\qty{288}{h}/\pi)\arcsin(1.02895/26.94) = \qty{3.50}{h}$。

\textbf{11.} $\sqrt{(1 + p)^2 - b^2} = (a/R_\star)\sin(\pi T_{14}/P) = 26.94\times
\sin(\pi\times3.23/288) = 0.949$（取 $\sin i \approx 1$）；$b^2 =
1.0587 - 0.9008 = 0.158$，$b = 0.40$。

\textbf{12.} $\cos i = b\,R_\star/a = 0.40/26.94 = 0.0148$，$i = \ang{89.15}$。

\textbf{13.} $R_\star/a = 1/26.94 = 0.037$：27 次中的一次机会。

\textbf{14.} $\sin i = 0.99989$：真实质量为 7.99 个地球质量，与\omterm{def:b3:exoplanets:rv}{最小质量}的
差距在 \num{1e-4} 以内。

\textbf{15.} $\bar\rho = \num{4.77e25}/\bigl(\tfrac43\pi(\num{1.531e7})^3\bigr) =
\num{4.77e25}/\num{1.502e22} = \qty{3180}{kg.m^{-3}}$，
\qty{3.18}{g.cm^{-3}}。

\textbf{16.} 该行星比质量相同的岩质行星大 $2.40/1.74 = 1.38$ 倍，
体积则为 2.6 倍。仅凭岩石无法形成它：它需要
一层轻质气体包层，氢与氦约占其质量的百分之几，
或是占很大比例的水。

\textbf{17.} 半径为 2.40 个地球半径，位于\omterm{def:b3:exoplanets:radius-valley}{半径谷}之上，落在 2--3 个地球半径
的峰区中，属于\omterm{def:b3:exoplanets:classes}{亚海王星}。

\textbf{18.} $g = Gm_p/R_p^2 = \num{6.674e-11}\times\num{4.77e25}/(\num{1.531e7})^2 =
\qty{13.6}{m.s^{-2}}$。

\textbf{19.} $L = 4\pi R_\star^2\sigma T^4 = 4\pi(\num{5.287e8})^2\times\num{5.670e-8}
\times5000^4 = \qty{1.245e26}{W}$，为 0.325 个太阳光度。

\textbf{20.} $S/S_\oplus = 0.325/0.0952^2 = 35.9$：远远位于
宜居带内缘之内。

\textbf{21.} $L/4\pi a^2 = \qty{4.88e4}{W.m^{-2}}$；$T_\mathrm{eq} =
\bigl(\num{4.88e4}\times0.7/(4\sigma)\bigr)^{1/4} = \qty{623}{K}$。

\textbf{22.} $H = kT/\bar m g = \num{1.381e-23}\times623/(2.3\times\num{1.661e-27}
\times13.6) = \qty{166}{km}$。

\textbf{23.} $2HR_p/R_\star^2 = 2\times\num{1.66e5}\times\num{1.531e7}/
(\num{5.287e8})^2 = \num{1.82e-5}$，即百万分之 18；高两个标高的一带
则为百万分之 36。

\textbf{24.} $H$ 缩小 $18/2.3 = 7.8$ 倍，为 \qty{21}{km}：每个标高
为百万分之 2.3，整条带为 4.6。若为百万分之 10，
富氢大气所显示的水汽带信号超过噪声的三倍以上；
蒸汽大气则显得平坦，与有云大气或
无大气的行星相似。

\textbf{25.} \textbf{\qty[text-series-to-math=true]{3.2}{g/cm^3}}。
\end{solution}
